% RESEARCH_DRAFT_SHA256: 9de46333543ad907b581d301e2bf76542af1772a36c1c6de2c91b39a3ebc9525
% Layout and source-reference wrappers expanded for the web renderer.
\documentclass[11pt]{article}
\usepackage[margin=25mm]{geometry}
\usepackage{fontspec,unicode-math}
\setmainfont{Latin Modern Roman}
\setmathfont{Latin Modern Math}
\usepackage{amsmath,amssymb,mathtools,xurl,hyperref}
\hypersetup{colorlinks=true,urlcolor=blue,linkcolor=blue}
\setlength{\parindent}{0pt}
\setlength{\parskip}{5pt}
\setlength{\emergencystretch}{3em}
\newcommand{\sref}[2]{\hyperlink{source:#1:#2}{[S#2]}}
\newcommand{\cataloguescope}[1]{\paragraph{Source question.}#1}
\newcommand{\floor}[1]{\left\lfloor#1\right\rfloor}
\begin{document}

\section*{A fixed parity component: a zero-dimensional CI calculation with all probability points checked}
\noindent UnsolvedMath ID 20001199; AIM-COMPUTATION-0037. Research attempt dated 11 October 2026.
\paragraph{Result in scope.}This manuscript contains a source-specific attempt and complete proofs of the stated auxiliary calculations. It makes no new solution or novelty claim for the entire target.
\section{Original problem and preserved definition}
The following definition and source list are reproduced from the repository. Any earlier inspection dates and reported literature claims in this preserved text describe that source record, not additional verification by this attempt.
% BEGIN PRESERVED DEFINITION
\subsection{Definitions and mathematical statement}

\paragraph{Shared notation.}Standard number systems are $\mathbb N,\mathbb Z,\mathbb Q,\mathbb R,\mathbb C$. Specialized notation and index ranges follow the cited source and the definitions below.


For finite discrete variables, conditional-independence constraints can be expressed by minors of conditional probability tables. Their algebraic zero set depends on the chosen constraints and ambient coordinates. A zero-dimensional algebraic variety has finitely many points over an algebraic closure. A matroid rank function $r$ assigns nonnegative integer ranks to subsets; its rank conditional-independence relation for disjoint $A,B,C$ is $r(A\cup C)+r(B\cup C)=r(C)+r(A\cup B\cup C)$.
\paragraph{Record scope and qualifications.}
The catalogue title is retained as metadata. Named certificates, contracts, bounds or frameworks in that title are not asserted as results of the original source extract.
The preceding definition of independence variety, including possible functional-independence constraints, is not supplied in this excerpt. The conjecture is retained without selecting an invented class or assuming that polynomial equations alone give the intended probability model. The supplied catalogue marks this record partially solved; its present research status was not independently checked.
The source text below is rendered literally where necessary; damaged formulas and composite extracts are not silently repaired into a new mathematical claim.
\paragraph{Statement recorded in the supplied source.}
Conjecture 23. When P belongs to such a zero dimensional independence variety then P

is the distribution of a random vector with its CI-structure equal to the CI-structure of a matroid.

Donald Richards

Donald Richards spoke on the problem of determining the number of solutions to generic maximum likelihood problems and the use of computational algebra, homotopy con-tinuation, and mixed volume in these calculations. \href{https://aimath.org/WWN/compalgstat/compalgstat.pdf}{[S2]}
\paragraph{Source status.}
The supplied catalogue labels this source record partially solved; that metadata was checked against the supplied snapshot on 2026-10-11. The label is preserved for provenance and is not an independent verification of present open status or mathematical correctness.
\subsection{Short English statement}
The source proposes a link between zero-dimensional independence varieties and matroid conditional independence.
\subsection{Sources}
\begin{itemize}
\item[S1] Ulam AI and the UnsolvedMath Contributors, UnsolvedMath: A Curated Collection of Open Mathematics Problems, record 20001199 (AIM-COMPUTATION-0037). Catalogue curation and metadata: CC BY 4.0. The source extract is retained; definitions, English summaries and qualifications are editorial additions.
\url{https://huggingface.co/datasets/ulamai/UnsolvedMath}.
\item[S2] AIM Workshop Problem Lists, source linked by catalogue record AIM-COMPUTATION-0037.
\url{https://aimath.org/WWN/compalgstat/compalgstat.pdf}.
\end{itemize}
% END PRESERVED DEFINITION
\subsection{Verbatim supplied catalogue question}
Conjecture 23. When P belongs to such a zero dimensional independence variety then P

is the distribution of a random vector with its CI-structure equal to the CI-structure of a matroid.

Donald Richards

Donald Richards spoke on the problem of determining the number of solutions to generic maximum likelihood problems and the use of computational algebra, homotopy con-tinuation, and mixed volume in these calculations.

\section{Why the undefined ambient class matters}
The source explicitly makes its conjecture conditional on independence
varieties being properly defined, including functional relations. With
only ordinary disjoint-set CI equations on finite variables, every
product distribution satisfies every such equation. If at least one
variable has two levels, fix all other marginals and vary that variable's
probability $t$. This gives a nonconstant affine line of normalized
complex product tables, and its positive real interval lies in the
probability simplex. Hence an ordinary CI locus of this unrestricted
kind cannot be zero-dimensional. A support or functional restriction
is essential to the intended route. The catalogue's parity title is
therefore treated as a test case, not as the general original target.

\paragraph{A fully specified finite calculation.}
Take binary $X,Y,Z$ with the fixed functional law
$Z=X\mathbin{\mathrm{xor}}Y$. In cell order $000,011,101,110$, let
the probabilities be $a,b,c,d$, with $a+b+c+d=1$ and all other cells
zero. Impose pairwise independence of each two variables. Their three
$2$ by $2$ marginal determinants are, up to sign,
\[
 ab-cd=0,\qquad ac-bd=0,\qquad ad-bc=0.
\]
For example, the $(X,Y)$ table is
$\left(\begin{smallmatrix}a&b\\c&d\end{smallmatrix}\right)$,
giving $ad-bc$; the $(X,Z)$ table is
$\left(\begin{smallmatrix}a&b\\d&c\end{smallmatrix}\right)$,
and the $(Y,Z)$ table is
$\left(\begin{smallmatrix}a&c\\d&b\end{smallmatrix}\right)$.
These verify every equation directly.

We solve this normalized system over $\mathbb C$, not just over the
nonnegative reals. If one coordinate is zero, all pair products of the
remaining three vanish; this follows by substitution in the three
equations (the equations are symmetric in the four coordinates as a
set). Thus at most one coordinate is nonzero, and normalization forces
one of the four coordinate vertices.
If all are nonzero, dividing the first equation by the second gives
$b/c=c/b$, so $b=\epsilon c$ with $\epsilon\in\{1,-1\}$.
Equivalently, solving the three equations gives
\[
b=\eta a,\qquad c=\xi a,\qquad d=\eta\xi a,
\qquad \eta,\xi\in\{1,-1\}.
\]
To verify completeness, $ad=bc$ first gives $d=bc/a$;
substitution in $ab=cd$ gives $c^2=a^2$, and substitution in
$ac=bd$ gives $b^2=a^2$. The displayed signs then follow.
Normalization is
$a(1+\eta)(1+\xi)=1$, so both signs must be $+1$ and $a=1/4$.
The normalized complex zero set therefore consists exactly of the
four vertices and $(1/4,1/4,1/4,1/4)$. It is finite, hence
zero-dimensional as a set; no assertion about multiplicity or a
scheme being reduced is needed. All five points are probability laws.

\paragraph{Exact CI structures at all five points.}
At each vertex all three variables are constant, so every conditional
independence holds. This is the rank-zero matroid on three elements.
At the uniform parity point each singleton has entropy one bit,
each pair has entropy two bits, and the triple has entropy two bits.
More directly, each pair is independent and uniform; each singleton is
a nonconstant function of the other pair; and given any one variable,
the other two are related by equality or complementation and are
therefore dependent. These exhaust the nontrivial disjoint CI triples
on three variables. For the rank function $r(A)=\min(|A|,2)$, the
rank-CI difference is zero for two unconditional singletons, one for
a singleton against the other pair, and one for the two remaining
singletons given the third. Thus the exact ordinary CI structure is
that of $U_{2,3}$. The stated functional dependencies also match: any
two variables determine the third, and no singleton determines another.
Every probability point of this particular finite variety satisfies
the proposed matroid conclusion.

\section{Attempted generalization and audit}
The computation fixes one deterministic parity map. General functional
dependence constraints can have many support components; their
normalizations, complex points and probability points must be analyzed
separately. Nothing in the five-point calculation proves that every
zero-dimensional component for every possible system has a matroid CI
structure. The original class of varieties is also not fully defined in
the source. The primary AIM Conjecture 23 and its preceding qualification
were inspected. The audit checked complex signs, normalization, zero
coordinates, all exact CI relations and the difference between finitely
many nonnegative points and a genuinely finite complex zero set. All
calculations were by hand; no Groebner basis or formal algebra package
was run. This is an elementary test case with no novelty claim and no
resolution of the general conditional conjecture.

\section{Process and contributions}
Mathematics and model-assisted adversarial self-audit: OpenAI gpt-6-astra, ultra reasoning, as dispatched for this research batch; exact serving revision not exposed. Codex (GPT-6; exact serving revision not exposed) materially assisted with orchestration and format review. No human mathematical authorship, mathematical edits, or independent human proof checking is asserted.
Contribution prompt A was applied to the full definition and supplied source row, then prompt B was applied adversarially to the derived results. The item-specific audit above records the scope checks and rejected extrapolations. Local assembly and structural checks do not establish mathematical correctness. The source catalogue status is preserved in TOP\_PROBLEM and is separate from this attempt status.
\paragraph{Integration review.} Codex (GPT-6; exact serving revision not exposed) reviewed the written arguments and expanded layout and source-reference wrappers. Repository metadata and the production web renderer were checked using Python 3.13 and JavaScript with MathJax 3.2.2. This model-assisted review is not independent human proof verification; no PDF compilation or formal proof checking is asserted.

\section{Completion Estimate}
\noindent\textbf{COMPLETION: 0\%}
The proved material is a rederivation of known or auxiliary facts rather than a substantiated new advance on the original target. A fixed binary parity system is solved completely and agrees with the conjectured conclusion, but the general class of functional-independence varieties is unspecified and the universal conjecture is not proved.
\section{References}
The complete original Sources subsection above is retained, including catalogue attribution and licence notices. The additional list records primary sources actually consulted for this attempt on 11 October 2026; uncited source links reproduced above were not all rechecked.
\begin{itemize}
\item American Institute of Mathematics, Computational Algebraic Statistics, prepared by Seth Sullivant, version 27 March 2004; the original numbered source question and its adjacent scope were inspected on 11 October 2026. \url{https://aimath.org/WWN/compalgstat/compalgstat.pdf}.
\end{itemize}
\section{Final status}
\textbf{Unresolved in this attempt.} A fixed binary parity system is solved completely and agrees with the conjectured conclusion, but the general class of functional-independence varieties is unspecified and the universal conjecture is not proved.
\end{document}
